\section{Implementation and Experiments}
\label{sec:experiment}

\subsection{Tool Implementation}

CBMC is a bounded model checker for C programs~\cite{clarke_biere_raimi_zhu_bounded_model_checking_2001,cbmc_2004,cbmc_architecture}.
Its verification flow includes front-end translation, symbolic execution, and solver-based
checking. This makes CBMC a suitable implementation target for this work for two reasons.

First, the approach in this report and the benchmark cases used in the evaluation are written in C. Therefore, CBMC can be used directly as the back-end checker after the checking model has been constructed.

Second, CBMC provides clear implementation points for the proposed methodology. The initial
program representation can be built around \texttt{goto-cc}, selected ISR insertion locations
can be computed before final verification, and OSEK/VDX API semantics can be handled during
symbolic execution. Thus, the implementation does not replace the underlying bounded model
checking engine. It changes the checking model and the symbolic execution behaviour so that
the final verification step receives a model that better reflects the executions of OSEK/VDX
applications.

The CBMC verification flow used in this report is shown in \Cref{fig:baseline-cbmc}. Source code is translated by the front-end into an internal program representation, that representation is symbolically executed to construct a solver formula, and the solver determines whether the required property holds under the chosen bound.

\begin{figure}[H]
  \centering
  \begin{tikzpicture}[
  node distance=7mm,
  box/.style={
    draw,
    rounded corners=2pt,
    align=center,
    minimum width=34mm,
    minimum height=10mm,
    fill=gray!8
  },
  arrow/.style={-{Latex[length=2mm]},thick}
]
  \node[box] (src) {C sources\\and harness};
  \node[box,below=of src] (goto) {Front-end\\and \texttt{goto-cc}};
  \node[box,below=of goto] (model) {GOTO model};
  \node[box,below=of model] (symex) {\texttt{goto-symex}\\SSA equation};
  \node[box,below=of symex] (solver) {SAT/SMT\\solver};
  \node[box,below=of solver] (trace) {Result and\\counterexample};

  \draw[arrow] (src) -- (goto);
  \draw[arrow] (goto) -- (model);
  \draw[arrow] (model) -- (symex);
  \draw[arrow] (symex) -- (solver);
  \draw[arrow] (solver) -- (trace);
\end{tikzpicture}

  \caption{CBMC architecture used as the implementation target for the proposed approach.}
  \label{fig:baseline-cbmc}
\end{figure}

This flow is a strong starting point for C verification, but it is not sufficient by itself for OSEK/VDX applications. The limitation is not symbolic execution itself. The limitation is that the checking model must also represent OSEK/VDX scheduling behaviours and must avoid checking ISR calls at program locations that do not affect the execution path.

The current implementation realizes the methodology inside CBMC by extending the flow in four places.
The \texttt{goto-cc} extension builds the initial checking model while preserving the OSEK/VDX project context. 
The \texttt{interleaving-analysis} stage produces a manifest of selected ISR insertion locations. 
The \texttt{aib} component uses this manifest to build the checking model with guarded ISR calls only at the selected locations. 
Finally, \texttt{goto-symex} is extended with the \texttt{os-api} layer so that recognized OSEK/VDX APIs update task states, events, the ready queue, and the next running task during symbolic execution.

Table~\ref{tab:cbmc-implementation-map} maps these implementation choices to the original CBMC stages. The purpose of the table is to show which parts of the CBMC flow are extended and which part remains unchanged.

\begin{table}[H]
  \centering
  \small
  \setlength{\tabcolsep}{4pt}
  \renewcommand{\arraystretch}{1.15}
  \begin{tabularx}{\linewidth}{@{}
    >{\raggedright\arraybackslash}p{0.18\linewidth}
    >{\raggedright\arraybackslash}p{0.16\linewidth}
    >{\raggedright\arraybackslash}p{0.26\linewidth}
    >{\raggedright\arraybackslash}X@{}}
    \toprule
    CBMC stage & Component role & Implemented extension & Effect on the checking model \\
    \midrule
    Front-end and \texttt{goto-cc}
    & Translate C into CBMC's internal representation.
    & Read the OSEK/VDX configuration and preserves project context.
    & Produces the initial checking model used by the later stages. \\

    \specialrule{0.2pt}{2pt}{2pt}

    Interleaving analysis
    & Not present in the CBMC flow..
    & Run \texttt{interleaving-analysis} before final verification.
    & Produces a manifest that records selected ISR insertion locations. \\

    \specialrule{0.2pt}{2pt}{2pt}

    Model building with selected ISR calls
    & Not present in the CBMC flow..
    & Use \texttt{aib} with \texttt{lib-interleaving-block}.
    & Builds a checking model with ISR calls only at the locations recorded in the manifest. \\

    \specialrule{0.2pt}{2pt}{2pt}

    Symbolic execution and \texttt{goto-symex}
    & Build the solver formula by symbolic execution.
    & Use the \texttt{os-api} layer to handle recognized OSEK/VDX APIs.
    & Updates task states, events, the ready queue, and the next running task according to OSEK/VDX scheduling behaviours. \\

    \specialrule{0.2pt}{2pt}{2pt}

    Solver back end
    & Solve the bounded query and reconstruct counterexamples.
    & Unchanged.
    & The final solving step stays the same; the input model has already been refined before solving. \\
    \bottomrule
  \end{tabularx}
  \caption{Implementation map from CBMC stages to the proposed approach.}
  \label{tab:cbmc-implementation-map}
\end{table}

\subsection{Benchmark Corpus and Experimental Setup}

The experimental corpus is drawn from two public benchmark: i-CBMC~\cite{icbmc_site} and IntAbs~\cite{intabs_repo}.
The corpus contains 43 public benchmark cases: 22 i-CBMC cases and 21 IntAbs cases.

The experimental setup follows the methodology described in \Cref{sec:method}. For each public case, two versions are checked under the same entry function, checked properties, and unwind bound. The CBMC version represents ISR calls broadly, while the proposed version uses the manifest-guided model-building step.

The main difference is how ISR calls are represented in the checking model. The proposed flow first identifies the ISR entry functions and the program locations that matter for each case. It then builds the checking model with ISR calls only at the selected locations, where
they can affect the execution path. Other locations are skipped because they only add unnecessary interleavings.

\subsection{Experiments}

The evaluation addresses three questions:

\begin{enumerate}[leftmargin=1.5em]
\item Does the proposed flow remove spurious bugs caused by unnecessary interleavings in the CBMC version?
\item Does the proposed flow reveal ISR-related behaviours that are not represented by the CBMC version?
\item What time and memory cost is introduced by the additional model-building stages?
\end{enumerate}

To answer these questions, each public case is reported in one of five outcome classes. 
A case has the \emph{same verification outcome} when both flows end in the same verification class. 
A case \emph{CBMC spurious bug removed} when the CBMC reports failure but the proposed model succeeds and the case audit supports the proposed version as the more accurate checking model. 
A case \emph{new error exposed by the proposed model} when the CBMC succeeds but the improved model reaches an error trace that the public audit supports. 
An \emph{unresolved mismatch} is a case where the two versions disagree, but the available evidence is not sufficient to decide which checking model is more accurate.
Finally, a case is \emph{not directly comparable} when the comparison is blocked by a resource limit or when the proposed flow cannot produce usable selected ISR insertion locations.

The end-to-end time for the CBMC version includes compilation and verification. 
The end-to-end time for the proposed version includes interleaving analysis, model building, compilation, and verification. 
Per-case time values are reported using median measured phase times, while peak memory is reported from the verification phase.

\begin{table}[H]
  \centering
  \small
  \renewcommand{\arraystretch}{1.15}
  \begin{tabularx}{\linewidth}{@{}
>{\raggedright\arraybackslash}p{0.28\linewidth}
>{\centering\arraybackslash}p{0.10\linewidth}
>{\raggedright\arraybackslash}X@{}}
    \toprule
    Comparison outcome & Cases & Interpretation \\
    \midrule
    Same verification outcome
    & 18
    & Both flows reach the same verification result. \\

    CBMC spurious bug removed
    & 8
    & The CBMC version reports a failure, but the proposed version succeeds because it removes unnecessary ISR interleavings. \\

    Unresolved mismatch
    & 8
    & The two versions produce different results, but the available evidence is not sufficient to decide which checking model is more accurate. \\

    Not comparable
    & 9
    & The comparison cannot be made directly because of resource limits or because no usable selected ISR insertion locations are produced. \\
    \bottomrule
  \end{tabularx}
  \caption{Outcome classes across the 43 public measured cases.}
  \label{tab:overall-outcome-summary}
\end{table}

\Cref{tab:overall-outcome-summary} shows that the proposed approach is mainly useful for building a more accurate checking model. At the same time, by reducing unnecessary interleavings, it can also substantially reduce the checking cost in many cases. The strongest positive result is the group of eight cases where CBMC spurious bugs are removed. In parcular, \textit{logger-oriented} case from the i-CBMC, the CBMC version reports \emph{VERIFICATION FAILED} with a median end-to-end time of 54,984 ms and a median verification-phase peak RSS of 297.4 MB. The proposed flow reports \emph{VERIFICATION SUCCESSFUL} for the same bound in 2\,025 ms and 16.1 MB. The main point is not only that the proposed version has a lower checking cost in this case. The more important point is that the CBMC flow permits an unnecessary ISR interruption point, which leads to a spurious verification failure. The proposed version checks the ISR entry as one atomic region, so this unnecessary interleaving is removed. As a result, the failure reported by the CBMC version is classified as a spurious bug.

The table also separates cases that should not be counted as clear wins or losses. Seven cases remain unresolved because the available evidence is not sufficient to decide which checking model is more accurate. Nine cases are not directly comparable because the comparison is blocked by resource limits or by missing selected ISR insertion locations. These cases are reported separately as limitations of the current evaluation, rather than being treated as artificial wins for either version. This makes the experimental summary more conservative and easier to interpret.

\subsection{Discussion}

The main cost of the proposed flow is the additional model-building work performed before final verification. On the 18 public cases where both versions produce the same verification result, the proposed flow is slower in 11 cases. Across this same-result subset, the median end-to-end time increases from 1,164 ms for the CBMC version to 3,380 ms for the proposed version. This result is expected because the proposed flow adds interleaving analysis and model-building stages before the final verification step.

However, this cost is not uniform across all cases. In some benchmark cases, the proposed version performs better than the CBMC flow in terms of checking time or memory usage. This mainly happens when the original program is larger or when the CBMC version explores many unnecessary ISR interleavings. In such cases, the cost of building a more restricted checking model can be compensated by the reduction in behaviours explored during final verification.

For the purpose of this report, this trade-off is reasonable. The objective is not to guarantee faster verification in every case. The first objective is to build a checking model that better reflects OSEK/VDX execution, especially task states, event state, the ready queue, and ISR-related behaviours that can affect the checked execution. The second objective is to make verification more feasible for larger source programs. In such programs, a broad ISR model may create many unnecessary interleavings, which can lead to high memory usage or timeouts in the CBMC flow. By selecting only the ISR insertion locations that can affect the checked execution, the proposed approach can reduce the behaviours explored during final verification. Therefore, the proposed approach mainly aims to make the checking model more accurate. At the same time, it can make verification more feasible for larger source programs where the CBMC flow is limited by high memory usage or timeout caused by unnecessary interleavings.
